\begin{historical}
    The idea that the dimensionality of physical space might be constrained by necessity predates the formal development of modern physics. Gottfried Wilhelm Leibniz suggested in the \emph{Discourse on Metaphysics} (1686) that the actual world should be understood as the one \QSOpen{}simplest in hypotheses and richest in phenomena,\QSClose{} implicitly framing dimensionality as subject to selection principles. In the eighteenth century, Immanuel Kant proposed that Newton's inverse-square laws implied the three-dimensionality of space, although this causal inference would later be reversed—the inverse-square law follows from spatial geometry via Gauss's theorem, not the reverse.
    
    A rigorous analytical approach began with Paul Ehrenfest's seminal 1917 paper \QSOpen{}In what way does it become manifest in the fundamental laws of physics that space has three dimensions?\QSClose{} He demonstrated that classical orbit stability requires exactly three spatial dimensions. In higher spatial dimensions ($N > 3$), the effective gravitational potential falls off too rapidly to maintain closed, bounded orbits; in lower dimensions, the dynamics become pathologically confined. Separately, the Huygens principle for the wave equation holds only in odd spatial dimensions $n \ge 3$, so $(3,1)$ spacetime is the lowest-dimensional case with sharp wavefronts and no interior tails.
    
    Gerald Whitrow's influential 1955 paper \QENOpen{}Why Physical Space has Three Dimensions\QENClose{} marked a turning point by explicitly connecting dimensional constraints to the possibility of life and observers. He argued that intelligent life capable of formulating physics could only arise in three spatial dimensions—an example of anthropic reasoning. Whitrow systematically examined how communication, neural networks, and information processing would fail in spaces of different dimensionality, establishing that the question \QSOpen{}why three dimensions?\QSClose{} might be answered by \QSOpen{}because otherwise we wouldn't be here to ask.\QSClose{}
    
    Freeman Dyson and Andrew Lenard's 1967 theorems on the stability of matter established that systems of electrons and nuclei interacting via Coulomb forces are stable of the second kind in three spatial dimensions. In higher spatial dimensions, Coulomb interactions scale differently and can lead to instabilities; in two or one spatial dimensions the Coulomb law changes and the stability analysis requires separate arguments.
    
    \end{historical}